\subsection{Controls: Support Mismatch, Sampling Geometry and Smoothing}
\label{sec:support_controls}

The rotated tile and the axis-aligned crop overlap over only $81\%$ of
their equal areas (Section~\ref{sec:utm_crop}). The main control keeps
that mismatch and takes the scene content out of the comparison. It is an
analytic random field, a sum of plane waves with power spectrum
$P(k)\propto k^{-\beta}$ multiplied by a Gaussian instrument response of
width $\sigma$, evaluated exactly at the level-20 cell centres and at the
native pixel centres of the same region. Both sides see the same
function, so whatever ratio survives comes from the sampling geometry
alone. Both parameters are measured: fitting
$\log P(k)=c-\beta\log k-k^2\sigma^2$ to the radially averaged spectrum of
each crop gives a median slope of $1.9$ and a median response of
\SI{11.9}{m} FWHM, against a nominal Sentinel-2 B04 response near
\SI{12.5}{m}, and natural images have exponents near
two~\cite{tolhurst1992amplitude,vanderschaaf1996powerspectra}. The same
fit places a median of $0.21\%$ of each crop's variance above half the
native Nyquist frequency.

Table~\ref{tab:support_controls} collects the controls, all
overhead-free. The synthetic rows evaluate the analytic field on both
lattices. The smoothed-native rows read the crop through the same
Clough--Tocher operator at points shifted by a fixed phase, at the native
density, and compress that raster like the crop: the native side gets
the \term{carrier}'s smoothing and nothing else. The common-support and
shifted-crop rows, which address the mismatch on subsets, are described
in Appendix~\ref{app:support}.

\begin{table*}[t]
\centering
\caption{Controls on the HEALPix/native size ratio at $\gamma=1$: median over the regions that bracket the target (their number in parentheses)}
\label{tab:support_controls}
\footnotesize
\begin{tabular}{@{}lrrrrr@{}}
\toprule
Control & $10^{-4}$ & $3\times10^{-4}$ & $10^{-3}$ & $3\times10^{-3}$ & $10^{-2}$ \\
\midrule
Imagery: Clough--Tocher tile / native & 2.07 (40) & 1.95 (40) & 1.58 (39) & 1.23 (37) & 0.93 (32) \\
Synthetic field, direct & 2.17 (40) & 2.08 (40) & 1.89 (39) & 1.46 (38) & 1.16 (32) \\
Synthetic field, via C--T & 2.17 (40) & 2.10 (40) & 1.90 (39) & 1.39 (38) & 0.94 (32) \\
HEALPix / smoothed native, phase 0.5 px & 2.22 (40) & 2.10 (40) & 1.83 (39) & 1.48 (37) & 1.24 (32) \\
HEALPix / smoothed native, phase 0.25 px & 2.19 (40) & 2.06 (40) & 1.78 (39) & 1.36 (37) & 1.10 (32) \\
Smoothed native / native, phase 0.5 px & 0.94 (40) & 0.92 (40) & 0.89 (40) & 0.84 (39) & 0.78 (35) \\
Smoothed native / native, phase 0.25 px & 0.96 (40) & 0.94 (40) & 0.92 (40) & 0.89 (39) & 0.86 (35) \\
Same ground, filled outside it & 2.08 (40) & 1.91 (40) & 1.65 (40) & 1.28 (36) & 0.97 (32) \\
Shifted crops, spread of the ratio & 0.012 (40) & 0.015 (40) & 0.022 (40) & 0.040 (40) & 0.057 (40) \\
\bottomrule
\end{tabular}
\end{table*}


With no content difference between the two supports, the direct ratio
falls monotonically from $2.17$ at $10^{-4}$ to $1.16$ at $10^{-2}$,
against $2.07$ to $0.93$ on the imagery. The control reproduces the
decline of the gap and its approximate magnitude, and lies above the
measured ratio at all five \term{targets}, by $5\%$ to $24\%$, more as
the \term{target} loosens, where its own sweep is partly censored
($n=38$ and $32$). Under $\gamma_{\mathrm{tile}}=1.30$ it reaches $0.89$
at the loosest \term{target} where the imagery reaches $0.72$. The rows
marked via C--T carry the same field to the cells through the paper's
own \term{transfer}, whose error against the exact field is $8.8\%$ of
the scene's standard deviation. The \term{transfer} changes nothing at
the three tightest \term{targets}, where the codec keeps every sample; at
the loosest it removes $18\%$ of the ratio, from $1.16$ to $0.94$,
against $0.93$ measured. The crossover appears in the control as soon as
the control contains the \term{transfer}, and the transferred control
tracks the measurement to within $1$--$20\%$.

The smoothed-native rows confirm this from the other side. They are a
strong-smoothing control at a fixed phase, not a matched experiment:
the \term{carrier} samples a different lattice at every phase. At half a
pixel, the phase at which the interpolant departs most from the samples,
the native raster compresses to $0.94$ of its size at the tightest
\term{target} and $0.78$ at the loosest (median over $40$ and $35$
regions), a reduction of $6\%$ to $22\%$ at equal error; at a quarter of
a pixel, to $0.96$ and $0.86$. Given that smoothing, the native side is
smaller than the \term{carrier} at every \term{target} and at both
phases. HEALPix over smoothed native falls from $2.22$ to $1.24$ at half
a pixel, with $32$ regions at the loosest \term{target} and the paired
test rejecting parity there ($p<10^{-4}$), and from $2.19$ to $1.10$ at
a quarter of a pixel ($p=4\times10^{-3}$). The half-pixel values sit
within $7\%$ of the direct synthetic control ($2.17$ to $1.16$) at every
\term{target}: once the native side carries the \term{transfer}'s
smoothing, what remains of the \term{carrier}'s cost is what sampling
the same field at $2.58$ times the density predicts. Across the original
comparison and the two phases, the loose-\term{target} ratio spans
$0.93$--$1.24$, with the two measured phases at $1.10$ and $1.24$. The
smoothing of the \term{transfer} is therefore enough to remove the
crossover. The controls on subsets (Appendix~\ref{app:support}) move the
ratio by $1\%$ of the effect at the tightest \term{target}; over
footprints that overlap by $74$--$83\%$, the support mismatch is unlikely
to explain the decline of the \term{size ratio}, though its contribution
to the narrow crossover at the loosest \term{target} is not isolated.
